% ITC2007 Course Timetabling Expert System — Proposal (LaTeX)
% Compile (recommended):
%   pdflatex proposal.tex
%   bibtex proposal
%   pdflatex proposal.tex
%   pdflatex proposal.tex

\documentclass[11pt,letterpaper]{article}

\usepackage[letterpaper,margin=1in]{geometry}
\usepackage[T1]{fontenc}
\usepackage{lmodern}
\usepackage{microtype}
\usepackage{hyperref}
\usepackage{booktabs}
\usepackage{enumitem}
\usepackage{xcolor}

\hypersetup{
  colorlinks=true,
  linkcolor=blue,
  urlcolor=blue,
  citecolor=blue
}

\title{A Rule-Based Prolog Expert System for Course Timetabling}
\author{\textit{Qiang Xue - 40300671, Aamna Fawad - 40329487, Avneet Kaur - 40326551} \\\\
Course: \textit{COMP 6591} \quad\quad Date: \textit{15 February 2026}}
\date{}

\begin{document}
\maketitle

\begin{abstract}
Course timetabling is a combinatorial optimization problem with hard feasibility constraints (e.g., no conflicts) and soft quality constraints (e.g., compact curricula, minimum working days)~\cite{bettinelli2015overview}. This project proposes a rule-based expert system implemented in SWI-Prolog to solve the ITC2007 Track~2 Curriculum-Based Course Timetabling benchmark~\cite{itc2007}. We will model the domain as a knowledge base of declarative constraints and build an inference/solving engine that constructs feasible timetables using a construction heuristic with backtracking and randomized restarts. The system will be evaluated on multiple ITC2007 instances using feasibility rate, penalty score, and runtime.
\end{abstract}

\section{Problem Statement}
We aim to build an expert system that generates weekly course timetables from ITC2007 Track~2 (Curriculum-Based Course Timetabling) instances~\cite{itc2007,de2022algorithm}. Each instance defines:
\begin{itemize}[nosep]
  \item a set of courses (teacher, number of lectures, minimum working days, number of students),
  \item a set of rooms (capacity),
  \item curricula (groups of courses that share students),
  \item and unavailability constraints (course cannot be scheduled in a specific time slot).
\end{itemize}
The goal is to assign every lecture to a \textit{day}, \textit{period}, and \textit{room} satisfying all hard constraints and minimizing a soft-constraint penalty.

\section{Objectives and Success Criteria}
\subsection{Objectives}
\begin{enumerate}[nosep]
  \item Parse ITC2007 \texttt{.ctt} data files and build an internal model.
  \item Encode hard constraints as Prolog rules and validate feasibility.
  \item Encode soft constraints and compute the objective (penalty).
  \item Implement a solver that produces feasible solutions.
  \item Provide reproducible experiments (scripts + CSV summaries) and write-up.
\end{enumerate}

\subsection{Success Criteria}
\begin{itemize}[nosep]
  \item \textbf{Correctness:} solutions pass a hard-constraint validator.
  \item \textbf{Coverage:} run on all provided instances (e.g., 21 competition instances).
  \item \textbf{Quality:} penalty is reported and shows improvement over a simple baseline (when feasible).
  \item \textbf{Reproducibility:} results re-generated with a single command producing a summary CSV.
\end{itemize}

\section{Dataset}
Instances are obtained from the official ITC2007 website:
\begin{center}
\url{https://www.eeecs.qub.ac.uk/itc2007/Login/SecretPage.php}
\end{center}

\section{Constraints (Domain Knowledge Base)}
\subsection{Hard Constraints}
Hard constraints must always hold:
\begin{itemize}[nosep]
  \item \textbf{Lecture completeness:} all required lectures for each course are scheduled.
  \item \textbf{Room conflict:} at most one lecture per room per time slot.
  \item \textbf{Teacher conflict:} a teacher teaches at most one lecture per time slot.
  \item \textbf{Curriculum conflict:} courses in the same curriculum do not overlap.
  \item \textbf{Unavailability:} forbidden (course, day, period) slots are never used.
\end{itemize}

\subsection{Soft Constraints (Penalty)}
Soft constraints define solution quality (lower is better). We will implement the Track~2 penalties:
\begin{itemize}[nosep]
  \item \textbf{Room capacity:} penalty when students exceed room capacity.
  \item \textbf{Minimum working days:} penalty when a course is spread across fewer than required days.
  \item \textbf{Curriculum compactness:} penalty for isolated lectures in a curriculum.
  \item \textbf{Room stability:} penalty when a course uses multiple rooms.
\end{itemize}

\section{Methodology and System Design}
\subsection{Expert-System Architecture}
We separate \textit{knowledge} (constraints) from \textit{inference} (solving)~\cite{guyette1994rule}:
\begin{itemize}[nosep]
  \item \textbf{Parser/Model:} reads \texttt{.ctt} into structured facts.
  \item \textbf{Rule base:} Prolog predicates representing constraints.
  \item \textbf{Solver engine:} constructs a schedule using guided search.
  \item \textbf{Validator/Scorer:} checks hard constraints and computes penalty.
  \item \textbf{Output:} writes \texttt{.sol} and CSV summaries.
\end{itemize}

\subsection{Solving Strategy (Construction-Only)}
Given the Phase~1 time constraints, we will focus on a construction-only solver that aims to reliably produce feasible timetables and report their penalties.
\begin{enumerate}[nosep]
  \item \textbf{Lecture ordering:} schedule the most constrained lectures first (e.g., high degree in conflicts, many unavailable slots).
  \item \textbf{Feasible placement:} choose a (day, period, room) assignment that satisfies all hard constraints, preferring placements that reduce immediate soft penalties (e.g., adequate room capacity).
  \item \textbf{Backtracking + randomized restarts:} if the solver gets stuck, it backtracks and/or restarts with different tie-breaking choices to improve feasibility rates.
\end{enumerate}
We will incorporate per-instance time limits and record runtime.

\section{Evaluation Plan}
\subsection{Metrics}
\begin{itemize}[nosep]
  \item Feasibility rate (\% of instances producing valid schedules).
  \item Total penalty (objective value) on feasible instances.
  \item Runtime per instance (seconds).
\end{itemize}

\subsection{Experimental Procedure}
\begin{itemize}[nosep]
  \item Run a simple baseline (e.g., first-fit greedy without restarts) to establish a reference.
  \item Run the construction solver (ordering + penalty-aware placement + restarts) under identical time limits.
  \item Aggregate results into a CSV summary and include tables in the final report.
\end{itemize}

\section{Timeline}
\begin{center}
\begin{tabular}{@{}ll@{}}
\toprule
\textbf{Week} & \textbf{Milestones} \\
\midrule
1 & Parse \texttt{.ctt}; instance validator; mini end-to-end run \\
2 & Implement hard constraints; solution validator; CI tests \\
3 & Implement soft constraints; penalty computation; baseline benchmarks \\
4 & Construction solver refinements; full benchmark runs; result tables \\
5 & Write report/paper; prepare presentation and demo \\
\bottomrule
\end{tabular}
\end{center}

\section{Work Breakdown Structure (WBS) and Contributions}
Table~\ref{tab:wbs} summarizes the planned work packages and expected contributions. Tasks may be adjusted as implementation findings emerge, but ownership ensures accountability and parallel progress.

\begin{table}[h]
\centering
\small
\begin{tabular}{@{}p{0.10\linewidth}p{0.46\linewidth}p{0.20\linewidth}p{0.18\linewidth}@{}}
\toprule
\textbf{WBS} & \textbf{Work Package} & \textbf{Owner(s)} & \textbf{Outputs} \\
\midrule
1.0 & Project setup, repo structure, CI, coding conventions & Qiang, Aamna & Working repo + CI \\
1.1 & ITC2007 Track~2 parser + instance model + sanity checks & Qiang & Parser + unit tests \\
2.0 & Hard-constraint rule base (teacher/room/curriculum/unavailability/completeness) & Aamna & Validator + tests \\
2.1 & Soft-constraint scoring (capacity, min working days, compactness, stability) & Avneet & Penalty module + tests \\
3.0 & Construction heuristic (greedy + restarts) + feasibility diagnostics & Qiang & Baseline solver \\
3.1 & Penalty-aware construction (ordering + room/day heuristics) & Avneet, Aamna & Construction solver (quality-focused) \\
4.0 & Benchmarking harness (batch runs, CSV summary, reproducibility) & Qiang & summary.csv \\
5.0 & Proposal + report + paper drafts (writing + figures/tables) & All & LaTex Documents  \\
6.0 & Presentation slides + demo script + rehearsal & All & Slides + demo \\
\bottomrule
\end{tabular}
\caption{Work Breakdown Structure and contribution plan.}
\label{tab:wbs}
\end{table}

% \section{Phase 1 Evaluation-Criteria Alignment}
% This proposal is structured to match the Phase~1 evaluation criteria:
% \begin{itemize}[nosep]
%   \item \textbf{Problem Definition (15\%):} Sections \textit{Problem Statement} and \textit{Constraints} define the ITC2007 Track~2 inputs, outputs, and constraints.
%   \item \textbf{Outcome Definition (15\%):} Section \textit{Objectives and Success Criteria} specifies measurable outcomes (validator pass, feasibility rate, penalty, runtime, reproducibility).
%   \item \textbf{Clarity (20\%):} Sections \textit{Methodology and System Design} and \textit{Evaluation Plan} describe components, data flow, and metrics.
%   \item \textbf{Completeness (20\%):} The proposal covers dataset, constraints, method, evaluation, timeline, and deliverables end-to-end.
%   \item \textbf{Conformance (15\%):} Uses the official ITC2007 Track~2 format, produces \texttt{.sol} outputs, and uses a validator/scorer consistent with the benchmark.
%   \item \textbf{Credibility and Quality of Literature (15\%):} Related Work grounds the approach in ITC2007 and established timetabling literature~\cite{itc2007,bettinelli2015overview,qu2009survey}.
% \end{itemize}

% \section{Risks and Mitigations}
% \begin{center}
% \begin{tabular}{@{}p{0.33\linewidth}p{0.6\linewidth}@{}}
% \toprule
% \textbf{Risk} & \textbf{Mitigation} \\
% \midrule
% Parser mismatches dataset format & Strict parsing using header counts; unit tests on real instances \\
% Feasibility hard to reach on some instances & Random restarts; better task ordering; add repair moves \\
% Optimization stuck in local minima & Randomized restarts; diversify ordering/tie-breaking; time-based stopping \\
% Runtime too high & Incremental constraint checks; profiling; tune restart/time limits \\
% \bottomrule
% \end{tabular}
% \end{center}

\section{Deliverables}
\begin{itemize}[nosep]
  \item Source code (Prolog), tests, and scripts for reproducible runs.
  % \item Benchmark CSV summaries and selected solution files.
  \item Proposal (this document), final report, short research paper, and presentation slides.
\end{itemize}

\section{Related Work (Starter)}
Curriculum-based course timetabling has been extensively studied in the literature~\cite{bettinelli2015overview}. The ITC2007 competition~\cite{itc2007} established standardized benchmarks and evaluation frameworks for this problem. Various approaches have been explored, including constructive heuristics and broader metaheuristic strategies~\cite{qu2009survey}. Recent work has focused on algorithm selection and instance space analysis to understand problem difficulty~\cite{de2022algorithm}. Rule-based expert system approaches have also been applied to university scheduling problems~\cite{guyette1994rule}, which motivates our Prolog-based approach.

\bibliographystyle{plain}
\bibliography{refs}

\end{document}
